\chapter{北京手记} % Introduction chapter suppressed from the table of contents


北京是全国的文化之都，来到北京后，晚上有空时我便去欣赏各种表演节目。连续三晚，我听了一场音乐会，看了两部话剧。

\hypertarget{ux56fdux5bb6ux5927ux5267ux9662ux7ba1ux5f26ux4e50ux56e2ux6f14ux594fux9a6cux52d2ux7b2cux4e09ux4ea4ux54cdux66f2}{%
\subsubsection{国家大剧院管弦乐团演奏马勒第三交响曲}\label{ux56fdux5bb6ux5927ux5267ux9662ux7ba1ux5f26ux4e50ux56e2ux6f14ux594fux9a6cux52d2ux7b2cux4e09ux4ea4ux54cdux66f2}}

到第六乐章的最后五分钟，两位定音鼓手齐声敲出``咚、咚、咚、咚''，仿佛为步兵擂鼓（又像赛龙舟的鼓手），与台上百余位乐手配合无间------包括6位打击乐手、40位铜管和木管乐手。每位乐手都全神贯注地盯着台前挥动的指挥棒，没有丝毫松懈，将全场听众带入最后的高潮。\\
在大剧院音乐厅，几百名听众已聆听了近100分钟的演奏，经历了第四乐章女中音的独唱，以及第五乐章合唱团与女中音的合唱。乐章之间没有片刻休息。\\
回想上一次的类似体验，是在香港大学音乐厅欣赏贝多芬第29钢琴奏鸣曲（``Hammerklavier''）的演出，全曲长达45分钟。钢琴家演奏结束时，全场听众与演奏者仿佛刚跑完一场马拉松，彼此都想互相拥抱，庆祝共同完成了这场非凡之旅。\\
论家中拥有多么昂贵高端的HiFi音响设备和大屏投影，都无法替代现场音乐会的感受。\\
非常认同美国中提琴家Kim
Kashkashian接受YouTube采访时所说的：``古典音乐的未来取决于现场听众与演奏者共同参与的音乐会（而不仅仅靠录制唱片、视频或DVD）。''\\
正如话剧一样，这种互动才赋予了表演生命力。\\
马勒的第一交响曲大获成功，为他后续交响曲的创作奠定了蓝图：

\begin{itemize}
\tightlist
\item
  使用超大规模的乐队
\item
  以歌曲作为主题
\item
  摒弃传统德国交响曲的格式，例如奏鸣曲式（sonata form）\\
\end{itemize}

他希望通过第三交响曲实现自己的创作理想。在奥地利萨尔茨堡（Salzburg）郊外阿特湖畔度假时，他终于从周围的高山湖泊、鲜花巨木和万物生灵中汲取了创作灵感。\\
最初，6个乐章都带有标题：

\begin{itemize}
\tightlist
\item
  牧神苏醒，夏季昂头阔步地迈进
\item
  草地上的花儿告诉我
\item
  森林里的动物告诉我
\item
  人类告诉我
\item
  天使告诉我
\item
  爱告诉我\\
\end{itemize}

但在1902年首演前，他删去了所有标题：``希望听众直接从音乐中去感受，不受标题的干扰。''\\
马勒的交响曲创作印证了此前提到过的弗里茨（Fritz）创新理论：一切创作都源于一个崇高却尚未实现的理想，感受到差距（张力），从而激发计划与行动。\\

\hypertarget{ux8bddux5267ux8fdbux5165ux9ed1ux591cux7684ux6f2bux957fux65c5ux7a0b-long-days-journey-into-night}{%
\subsubsection{话剧：进入黑夜的漫长旅程 (Long Day's Journey into
Night)}\label{ux8bddux5267ux8fdbux5165ux9ed1ux591cux7684ux6f2bux957fux65c5ux7a0b-long-days-journey-into-night}}

话剧表演结束后，我一看手表，已经快10点了。不知不觉中连续看了将近150分钟，没有片刻休息，现在才觉得屁股有些酸痛------因为坐得太久了。\\
虽然全剧只有一家四口------父母和两个儿子------但通过夫妻间、父子间、兄弟间接连不断的对话，让观众从他们早晨的笑脸背后，渐渐感受到每个人的痛苦、无奈与绝望。尤其是父亲与小儿子Edmund在深夜交心时，父亲诉说自己小时候一家人从爱尔兰移民到美国的贫困经历，后来虽然赚了钱，却深知赚钱不易，结果被家人误解为``守财奴''，认为他不顾家。那一刻，我实在忍不住泪水，掏出手帕擦了擦脸。\\
这部戏是美国经典剧目，作者尤金·奥尼尔（Eugene
O'Neill）生前三次荣获普利策奖。这是他自认为最好的作品，在他死后的1956年出版，并在1957年再次获得普利策奖。\\
如果只是借用别人的故事，内容很难有深度，必须植根于个人真实经历。奥尼尔年轻时考入著名的普林斯顿大学，一年后却因酗酒被劝退。正因如此，他才能如此深刻地刻画《进入黑夜》中长子Jamie------一个酒鬼------的想法、心态以及酗酒的缘由。其实，他不过是借剧中四个角色，将自己的亲身经历传递给观众罢了。\\
田纳西·威廉斯（Tennessee Williams）的成名作《玻璃动物园》（The Glass
Menagerie）也是如此，写的就是自己的家庭：轻度残疾的妹妹找不到工作，母亲想尽办法为她寻找伴侣等。

\hypertarget{ux8bddux5267ux5b8cux7f8eux964cux751fux4eba}{%
\subsubsection{话剧：完美陌生人}\label{ux8bddux5267ux5b8cux7f8eux964cux751fux4eba}}

虽然票价最贵（300元），却是三晚中最平淡的演出。七位演员很专业也全情投入，但却未能打动我。这部剧改编自2016年的一部意大利电影，讲述了七位老友（包括三对夫妇）晚上在其中一家聚餐的故事。晚餐时，他们玩了个游戏：将手机放在桌上，公开分享收到的消息和电话。结果像炸弹接连引爆------揭露了家庭矛盾、混乱的男女关系等。起初觉得有趣，但揭秘越多就越感觉麻木，总觉得这些故事像是导演凭空想象，缺乏真实感。我也不喜欢演员全程佩戴麦克风讲话，像在听晚间电台广播，少了与观众互动的氛围。

\hypertarget{ux521bux4f5cux6cc9ux6e90}{%
\subsection{创作泉源}\label{ux521bux4f5cux6cc9ux6e90}}

音乐、话剧等艺术形式为创作者和表演者提供了无限的创作空间。\\
从400年前莎士比亚五花八门的喜剧与悲剧，到250年前莫扎特的歌剧、交响曲、协奏曲中千变万化的主题变奏，无不让观众耳目一新。（近代披头士Beatles、Bob
Dylan的歌曲也同样深深吸引着我们。）因此，我特别喜欢买票去看现场演出，亲身感受创意带来的震撼。如果我们观察小孩，会发现他们在成长初期想象力丰富、充满创造力，为什么长大后反而缺失了呢？\\
``这是教育的后果。''美国宾夕法尼亚大学沃顿商学院（U of Penn, Wharton
Business School
）的Ackoff教授说，``主要原因是他们要应付学校里的考试和功课。''\\
在学校，学生面对试题时，第一反应是猜测老师想要的正确答案。学校和老师会不断``教育''孩子追求标准答案，拿满分。

\framebox{%
\begin{minipage}[t]{0.97\columnwidth}\raggedright
比如：3 + 5 =
？读过书的人会立刻回答8，而不会思考问题的场景。如果不是十进制，答案可能是10或12。\strut
\end{minipage}} 

只有通过多听、多读、多看，不断拓宽视野，才能提升个人的原创能力。但千万不要局限于自己专业领域的内容。\\
有人觉得奇怪，为什么香港大学的数学系归属文学院，而不是工程学院或理学院？虽然许多工程和科学发明都要依赖数学，但数学与哲学、音乐也并非毫无关联。著名数学家伊藤先生曾说：``数学其实能帮我们发现大自然的壮美与真理，比许多哲学理论更贴近现实。''所以，我们不应给自己设下框框。

\hypertarget{ux4e0dux65adux5b66ux4e60}{%
\subsubsection{不断学习}\label{ux4e0dux65adux5b66ux4e60}}

要让学生有学习的动机和动力，最佳的方式就是让学生当``老师''来指导学生。回顾我20年前教授PMP和项目管理，以及2018年讲授ACP敏捷课程时，为了准备教材、给学员上课、回答问题，我在这个过程中确实加深了自己对敏捷项目管理的理解。

学习可以分成以下阶段：

\begin{itemize}
\tightlist
\item
  数据 (Data)
\item
  信息 (Information)
\item
  知识 (Knowledge)
\item
  理解 (Understanding)
\item
  智慧 (Wisdom)
\end{itemize}

从理解数据信息，变为智慧。\\
培训过程中，与学生互动，解答他们的疑问，让我能仔细深入思考，从单是理解，提升到实际应用。\\
反之，有些大学生学编码，利用AI技术,搜索对应老师题目的答案，便交作业。(程序可能连编译都不通过)\\
大家不一定都有机会讲课，但都可以记录与分享经验教训。\\
这些真实经历，如果没有及时记录下来，后面就很容易忘记。假如我没有把刚刚连续看演出的感受记在纸上，过了一个月后肯定都变得模糊。所以我也尽力每天把过程记录下来。\\
若想不断学习，除了多听多看外，还要多思考、多写，因为如果没落到纸上都只是模糊的印象。大家有每天写日记的习惯吗？\\
现在科技发达，比120年前民国前要用毛笔写字轻松多了，但当时很多伟人还是每天不停写日记。我们现在可以利用语音识别等技术，只花15到30分钟便可以把每天经历的重点记下来。\\
提示：虽然每天记录、写草稿，估计最终能有十分之一留下来发布已经不错。

\framebox{%
\begin{minipage}[t]{0.97\columnwidth}\raggedright
在美国，越来多越多年轻女性参加以前只有男性的体育运动比赛，导致很多女性因缺乏准备、锻炼和热身，会腿部前十字韧带受损，做紧急手术。记者Cynthia
Gorney
就想按这个主题给纽约时报杂志写专题文章。整个过程像漏斗一样：先把各研究、访问得到的相关资料放进漏斗，到了交稿截点前开始整理资料，把零散的资料组织成有架构的文章，分章节、如何开头、如何结尾等。然后杂志编辑会按读者群体的关注点精简初稿，比如她初稿也包括研究过程中她觉得挺有兴趣的替换膝盖手术，她甚至特意飞到Cincinnati观察医生做手术，但编辑觉得这些跟主题无关，非读者关注，把这些部分都删除掉。经过多轮的过滤、修正，最后在杂志出版的文章4000字，共7页。（详见参考书）
\strut
\end{minipage}} 

但是确实要靠这去芜存精的过程才可以保证发布文章的可读性和吸引力。\\
前面说过乔布斯在大学演讲的时候给学生提3个故事，第一个故事分享当初在大学学美术字（Calligraphy）时只因为很感兴趣，没有细想有什么用，到后来在苹果公司为Macintosh个人电脑设计字体时，才发现能用上。他用这故事说明很多时候过后回顾才知道以前那些点可以串联起来，对后面有作用。我也希望能在日后能用上之前自己的听、看、读、写。

\href{文件:25A图1.jpg}{500px}

和上面话剧、音乐会一样，必须与其他人交流，不断完善才能出好东西。总起来说，越接触各领域大师的创作，就越觉得自己渺小，但只要知道自己是今个月比上个月有进步，比半年前有进步就可以了。\\
不是每个人都可以成为乔布斯或马勒、尤金•奥尼尔，但都可以成为本职工作的专业知识工作者。\\
我喜欢看演出的另一个原因是想感受他们的专业性：台上的音乐家、指挥、歌唱家、演员，和幕后的导演、戏剧作家、作曲家等。\\
专业和业余的区别：已经从多轮经验教训改善、精益求精，能让客户完全放心。\\
要专业，除了具备本职工作的知识技能，也需要有前瞻性。要不断改善，除了具备相关知识以外，更重要是态度：必须有自我不断完善的愿望(self-motivation)和个人规范(discipline),
只有这样做才可以保持自己的专业地位。\\
在当今竞争激烈、需求不断变化的大环境下，企业每位员工都要准备好自己，做专业XX师。员工要专业，也需要上面提到多听多读多看多写，打破现有的框框，全局了解客户的要求，动脑筋分析，不断改善，不然这些都只是梦想。

\hypertarget{ux53c2ux8003-reference}{%
\section{参考 Reference}\label{ux53c2ux8003-reference}}

\begin{enumerate}
\tightlist
\item
  Kramer. Mark and W. Call ''Telling True Stories. '' (Penguin , 2007)\\
\end{enumerate}


